This project addresses the needs within modern server infrastructure: continuous availability, security and responsiveness. It is very important to have services uninterrupted and reliable for both businesses and end-users. To that end various themes and works have been analyzed.

\subsection{Infrastructure as Code}
Infrastructure as Code (IaC) is the basis of how all modern projects function. It allows configuration to be automated, consistent and repeatable for managing available resources. Using IaC drastically increases the speed of development, deployment and testing of server infrastrcture \cite{Saxena2024DevOps, YUNGAICELANAULA202264} 

\subsection{Reactive vs Proactive security}
A major strength in security management is proactively looking for security vulnerabilities. Instead of merely responding to security incidents after they occur it is a lot more safer and productive to seek them out before they happen. This proactive stance reduces vulnerabilities and enhances overall security, aligning with cybersecurity best practices.

\subsection{Software Defined Networking}
Software Defined Networking (SDN) is a method of managing networks by using software instead of configuring hardware directly. It allows separating the network’s control plane, which decides how data should flow. This separation makes it easier to dynamically configure and manage the network. SDN allows for quicker adaptations to changing network conditions and needs.

\subsection{Zero-Touch Provisioning}
Zero-Touch Provisioning (ZTP) \cite{Watsen2019SZTP} automates the initial deployment and configuration of network devices without manual intervention. By leveraging predefined configurations and scripts, ZTP accelerates the setup process, reduces human error, and ensures consistency across network devices. This approach is particularly beneficial in large-scale deployments where manual configuration would be time-consuming and error-prone.

\subsection{Intent-Based Networking}
Intent-Based Networking (IBN) shifts network management from manual configurations to a model where administrators define high-level business intents\cite{Manias2024Intent}. The network then automatically translates these intents into configurations and enforces them across the infrastructure.

\subsection{Artificial Intelligence for IT Operations (AIOps)}
AIOps integrates artificial intelligence and machine learning into IT operations to automate the detection, diagnosis, and resolution of network issues\cite{diaz2023aiops}. By analyzing vast amounts of network data in real-time, AIOps can predict potential problems, optimize performance, and enhance security measures. This proactive approach minimizes downtime and improves overall network reliability.